\section{Transition Certificates}
\label{sec:method}
In this section, we propose a new type of proof certificates, called transition certificates, to verify that a system $\mathfrak{S}$ satisfies an LTL formula $\phi$. The core idea of transition certificates is to reduce the verification problem to analyzing the behavior of accepting transitions in the NBA $\mathcal{A}$ for $\neg\phi$. By the definition of the transition-based acceptance condition, a word is accepted by $\mathcal{A}$ if and only if accepting transitions occur infinitely often along one of its runs. Therefore, to prove $\mathfrak{S} \models_{\mathfrak{L}} \phi$, it suffices to show that for all trajectories of $\mathfrak{S}$, every accepting transition in $\mathcal{A}$ occurs at most finitely many times. We introduce two types of transition certificates: \emph{transition safety certificates} and \emph{transition persistence certificates}. The two certificates address the problem from fundamentally different angles and concern the set $Acc^{start}$, which comprises the starting states of all accepting transitions in $\mathcal{A}$ (cf. equation~\ref{atstart}).

\begin{itemize}
    \item A \emph{transition safety certificate} proves that the system's trajectories can never reach $Acc^{start}$, thereby guaranteeing that related accepting transitions \textbf{never occur}.
    \item A \emph{transition persistence certificate} allows the system to reach $Acc^{start}$ but proves that related accepting transitions can only be taken \textbf{finitely often}.
\end{itemize}


\subsection{Transition Safety Certificate}

The objective of a transition safety certificate is to prevent any accepting transition from occurring. To achieve this, we verify that accepting transitions never occur by showing that $Acc^{start}$ is unreachable in the product of the system and the automaton.

\begin{definition}[transition safety certificate]
  \label{transition-safety-certificate-definition}
  Given a system $\mathfrak{S} = (\mathcal{X}, \mathcal{X}_0, f)$, a labelling function $\mathfrak{L}:\mathcal{X} \rightarrow 2^{AP}$ and an NBA $\mathcal{A} = (2^{AP}, Q, q_0, \delta, Acc)$. Construct the product system $\mathfrak{S} \times \mathcal{A} = (\mathcal{Y}, \mathcal{Y}_0, F)$, where $\mathcal{Y} = \{(x, q) \,|\, x \in \mathcal{X}, q \in Q \}$, $\mathcal{Y}_0 = \{ (x_0, q_0) \,|\, x_0 \in \mathcal{X}_0 \}$, and $F(x,q) = \{(x^{\prime}, q') \,|\, x^{\prime} = f(x) \land (q, \mathfrak{L}(x), q') \in \delta\}$. For $q_k \in Acc^{start}$, define the unsafe set $\mathcal{Y}_u = \{(x, q_k) \,|\, x \in \mathcal{X}\}$. A bounded function $B_{k}(x, q): \mathcal{X} \times Q \rightarrow \mathbb{R}$ is a \emph{transition safety certificate} with respect to $q_k \in Acc^{start}$ if for all $(x,q) \in \mathcal{Y}$, $(x_0, q_0) \in \mathcal{Y}_0$, $(x^{\prime}, q^{\prime}) \in F(x, q)$, we have:
  \begin{align}
    &B_{k}(x_0, q_0) \geq 0, \label{btsc1} \\
    &B_{k}(x, q_k) < 0, \label{btsc2} \\
    &B_{k}(x, q) \geq 0 \Rightarrow B_{k}(x^{\prime}, q^{\prime}) \geq 0. \label{btsc3}
  \end{align}
\end{definition}

% Intuitively, these conditions ensure that: 1) the certificate value is non-negative at the initial state; 2) it becomes negative precisely on the unsafe set $\mathcal{Y}_u^k$ associated with $q_k$; and 3) once non-negative, its value remains non-negative along any trajectory of the product system. The certificate $B_k$ guarantees $q_k$ is unreachable, preventing all accepting transitions from $q_k$. The following lemma formalizes this guarantee. Its proof is provided in the Appendix~\ref{appendix-proof}.

Intuitively, condition (\ref{btsc1}) guarantees that the certificate value is non-negative at the initial state of the product system; condition (\ref{btsc2}) ensures that the certificate becomes negative precisely on the unsafe set $\mathcal{Y}_u^k$ associated with $q_k$, thereby marking states where accepting transitions could originate; condition (\ref{btsc3}) implies that once the certificate value is non-negative at a state, it remains non-negative along any trajectory of the product system, thus preserving the safety invariant.

Collectively, these conditions ensure that the certificate $B_k$ guarantees $q_k$ is unreachable, preventing all accepting transitions from $q_k$. The following lemma formalizes this guarantee. Its proof is provided in Appendix~\ref{appendix-proof}.

\begin{lemma}
  \label{transition-safety-lemma}
  If a transition safety certificate $B_k$ with respect to $q_k \in Acc^{start}$ can be found, then all runs of words in $\mathfrak{L}(\mathfrak{S})$ never visit state $q_k$, thereby preventing all accepting transitions starting from $q_k$.
\end{lemma}

The next theorem follows directly from applying the above lemma to every state in $Acc^{start}$. If a transition safety certificate can be found for each state in $Acc^{start}$, then no accepting transition can be taken, ensuring the LTL specification holds.

\begin{theorem}
  \label{transition-safety-LTL}
  Given a system $\mathfrak{S} = (\mathcal{X}, \mathcal{X}_0, f)$, a labelling function $\mathfrak{L}:\mathcal{X} \rightarrow 2^{AP}$ and an NBA $\mathcal{A} = (2^{AP}, Q, q_0, \delta, Acc)$ for $\neg\phi$. If for any $q_i \in Acc^{start}$, a transition safety certificate $B_i$ can be found, then it can be guaranteed that $\mathfrak{S} \models_{\mathfrak{L}} \phi$.
\end{theorem}

\begin{proof}
By Lemma~\ref{transition-safety-lemma}, all start states of accepting transitions are unreachable. Consequently, accepting transitions never occur, so the system language is not accepted by the automaton for $\neg\phi$. This proves $\mathfrak{S} \models_{\mathfrak{L}} \phi$.
\end{proof}

\subsection{Transition Persistence Certificate}
In contrast, the \emph{transition persistence certificate} offers an alternative approach. Instead of completely preventing the system from reaching states where accepting transitions may happen, this certificate proves that even if such states are reached, the corresponding accepting transitions can occur only a finite number of times. This is achieved by requiring a decrease in a certificate value each time an accepting transition is taken.

\begin{definition}[transition persistence certificate]
  \label{transition-persistence-certificate-definition}
  Given a system $\mathfrak{S} = (\mathcal{X}, \mathcal{X}_0, f)$, a labelling function $\mathfrak{L}:\mathcal{X} \rightarrow 2^{AP}$ and an NBA $\mathcal{A} = (2^{AP}, Q, q_0, \delta, Acc)$. For $(k, l) \in Acc^{pair}$, define the indicator function $I_{k,l}(x) = 1$ if $(q_k, \mathfrak{L}(x), q_l) \in Acc^{k,l}$, and $0$ otherwise. A bounded function $B_{k, l}(x): \mathcal{X} \rightarrow \mathbb{R}$ is a \emph{transition persistence certificate} with respect to $Acc^{k,l}$ if there exists $\epsilon > 0$ such that for all $x_0 \in \mathcal{X}_0$, $x \in \mathcal{X}$, $x^{\prime} = f(x)$, we have:
  \begin{align}
    &B_{k, l}(x_0) \geq 0, \label{btpc1} \\
    &B_{k, l}(x) \geq 0 \Rightarrow B_{k, l}(x^{\prime}) \geq 0, \label{btpc2} \\
    &B_{k, l}(x) \geq B_{k, l}(x^{\prime}) + I_{k,l}(x)\epsilon. \label{btpc3}
  \end{align}
\end{definition}

%Intuitively, 
%The certificate $B_{k,l}$ is non-negative initially (\ref{btpc1}), remains non-negative along trajectories (\ref{btpc2}), and decreases by $\epsilon$ whenever the state's label matches an accepting transition from $q_k$ to $q_l$ (\ref{btpc3}). 

Intuitively, condition (\ref{btpc1}) guarantees that the certificate $B_{k,l}$ is non-negative at the initial state of the system; condition (\ref{btpc2}) ensures that $B_{k,l}$ remains non-negative along any trajectory of the system; condition (\ref{btpc3}) implies that $B_{k,l}$ decreases by at least $\epsilon$ whenever the state's label matches an accepting transition from $q_k$ to $q_l$, and remains non-increasing otherwise.

% By (\ref{btpc1}) and (\ref{btpc2}), $B_{k,l}$ remains non-negative along any trajectory starting from $x_0 \in \mathcal{X}_0$. By (\ref{btpc3}), $B_{k,l}$ decreases by at least $\epsilon$ each time $I_{k,l}(x) = 1$. 

Since accepting transitions from $q_k$ to $q_l$ can only occur when $(q_k, \mathfrak{L}(x), q_l) \in Acc^{k,l}$, the number of such transitions is bounded by $B_{k,l}(x_0)/\epsilon$, ensuring finite occurrence. This yields the following lemma. Its proof is provided in the Appendix~\ref{appendix-proof}.

\begin{lemma}
  \label{transition-persistence-lemma}
  If a transition persistence certificate $B_{k,l}$ can be found, then all runs of words in $\mathfrak{L}(\mathfrak{S})$ make accepting transitions from $q_k$ to $q_l$ happen finitely often.
\end{lemma}

The next theorem is obtained by applying the lemma to every accepting transition pair.

\begin{theorem}
  \label{transition-persistence-LTL}
  Given a system $\mathfrak{S} = (\mathcal{X}, \mathcal{X}_0, f)$, a labelling function $\mathfrak{L}:\mathcal{X} \rightarrow 2^{AP}$ and an NBA $\mathcal{A} = (2^{AP}, Q, q_0, \delta, Acc)$ for $\neg\phi$. if for any $(k, l) \in Acc^{pair}$, a corresponding transition persistence certificate $B_{k,l}$ can be found, then it can be guaranteed that $\mathfrak{S} \models_{\mathfrak{L}} \phi$.
\end{theorem}

\begin{proof}
By Lemma~\ref{transition-persistence-lemma}, every accepting transition occurs only finitely often. The automaton for $\neg\phi$ therefore has no accepting run on any system trace. This proves $\mathfrak{S} \models_{\mathfrak{L}} \phi$.
\end{proof}

The two certificate types are complementary. A safety certificate is appropriate when every accepting-transition start state can be proved unreachable. This requirement can be too strict for persistence-style behavior, where the system may reach such states but still take accepting transitions only finitely often. A persistence certificate covers that case by proving a decrease whenever the accepting transition is taken. The converse limitation also matters. For $\mathbf{G}\neg a$ with the NBA for $\mathbf{F}a$, condition~(\ref{btpc3}) would require an infinite descent of a bounded value whenever the accepting self-loop is enabled. A persistence certificate cannot satisfy this obligation, while a safety certificate can still prove the property by showing that the accepting-transition start state is unreachable. The main theorem uses this complementarity. For each accepting-transition start state, it is enough to discharge the corresponding accepting transitions by either safety reasoning or persistence reasoning.


In summary, transition certificates give two proof mechanisms for excluding accepting runs of the automaton for $\neg\phi$. The construction keeps the certificate variables tied to transition starts or transition pairs, which reduces the search space relative to closure-certificate formulations used in our comparison.

